%% notes-tex/modeling/scaling-laws/sections/3-counting.tex

\section{Counting $N$, $D$ and $C$ when the examples are strains\secstatus{todo}}
\label{sec:counting}

The forms were fit on language models, where the units are settled: a parameter
is a weight, a datum is a token, and compute is $6ND$. For a model whose examples
are genotype, environment and phenotype records, each unit has to be chosen, and a
curve is only as defensible as the choice. The rule that resolves most of the
cases is that $D$ counts in the unit the loss averages over, and that the unit is
held fixed across every run in the sweep.

\notesec{$D$ is instances consumed, not dimensions}
$D$ is the number of training examples the optimizer has processed when the loss is
measured: the number of records in the training subset times the number of
epochs, in the same unit the loss is averaged over. For a language model that unit
is the token because the cross-entropy is per token; for a model trained on one
loss per record it is the record. The dimension of an example, meaning the number
of features in a record, the number of genes in the genome the model attends over,
or the length of a sequence embedding, is not $D$. Dimension changes what a record
is and how much compute each record costs, and it changes $N$ through the width of
the input layer, but a sweep over dimension is a sweep over architecture, not over
data, and belongs on a different axis of the ablation. Two consequences follow.

\begin{itemize}
\item A record with several supervised labels, a genotype measured on several
phenotypes, counts as one record under a summed loss and as one example per label
under a per-label loss. Either is a valid unit; what breaks the curve is changing
the unit between runs, or reporting a loss averaged over one unit against a $D$
counted in another.
\item A gene-token model that reads a genome as a set of tokens does not consume
$D \times 6{,}000$ tokens in the Kaplan sense: the loss is still one number per
record, and the tokens inside a record are the dimension, not the data. The
per-record compute does scale with that token count, which is why $C$ is measured
rather than computed for such a model (below).
\end{itemize}

\notesec{Repeated epochs count for less}
Language models are trained for about one epoch, so $D$ equals the corpus. A
genotype-to-phenotype model trains for many epochs over a fixed corpus, and a
repeated record is not worth a fresh one. Muennighoff et al.\
(arXiv:2305.16264)\external{arXiv:2305.16264, not in the mirror} fit the
Chinchilla form with an \term{effective data} term in which repeated tokens
decay in value: up to about four epochs the loss is close to what the same number
of fresh tokens would give, and beyond that the return on repetition falls off
quickly. The usable statement for a data-limited project is that $D$ in
Equation~\ref{eq:chinchilla} is the number of distinct records, and the epoch
count is a training hyperparameter held fixed across the sweep, not a multiplier
on $D$. A $D$-sweep is then a sweep over the size of the distinct training set,
which is what asks the question the project needs answered: whether a new
experiment or a bigger model will pay more.

\notesec{Subsample by the unit that leaks}
A $D$-sweep draws nested subsets of the training set, 1/64, 1/16, 1/4 and the
whole, against one fixed evaluation set. If test genotypes share genes or
perturbations with training genotypes, adding data partly measures memorization of
those shared parts, and the apparent exponent is inflated. The subsets are
therefore drawn by the unit that leaks: by gene, or by the query pair a screen was
built around, so that a smaller subset removes whole genes or whole screens rather
than thinning records uniformly. The 025 additive-baselines document showed the
size of this effect on the trigenic build, where the additive ridge falls from
0.400 test Pearson under a random-over-records split to 0.185 under a
query-pair-disjoint one; a $D$-sweep drawn by record on that build would scale
the first number, not the second.

\notesec{$N$: count what the loss responds to}
$N$ is the number of trainable parameters that the loss responds to as capacity.
Kaplan's exclusion of the embedding matrix has a direct analog: a gene-token model
carries an embedding table with one row per gene, whose size is set by the genome
rather than by the architecture, and that table should be reported separately
from the rest, with the curve fit on the rest. Sizes are then produced by scaling
width and depth together in a fixed ratio, five to seven of them spaced evenly in
$\log N$ over at least one and a half to two decades.

\notesec{$C$: measure it}
The $6ND$ rule is transformer-specific and per token. For a model whose per-record
cost depends on how many gene tokens a record carries, on graph message passing,
or on a sequence encoder run once per gene, $C$ is read from a profiler for one
training step at each size and multiplied by the number of steps.
\texttt{torch.utils.flop\_counter.FlopCounterMode} does this for a PyTorch model
without instrumentation. The IsoFLOP construction of panel~g then uses the measured
cost per record in place of $6N$.

\notesec{$E$ has a physical reading}
For a language model $E$ is the entropy of text and is not independently
measurable. For a phenotype measured with replicates it is. Under a mean-squared
loss the best possible predictor has expected error equal to the measurement-noise
variance, so the replicate variance of the evaluation set is a lower bound on $E$
in the loss's own units, and it is available before any model is trained. The
released standard-deviation columns of the Costanzo and Kuzmin screens, whose
replicate structure is recorded in
\file{notes/torchcell.datasets.scerevisiae.costanzo2016.noise-computation.md},
give the per-record $\sigma$ from which that bound is computed. A fitted $E$ that
falls below it says the fit is wrong, and one far above it says the model family
is leaving something on the table that neither more data nor more parameters will
recover. Which of those holds on the trigenic build has not been measured.
